\section*{Instruction Sheet 4: Supervised Modeling of a Radioactive Decay Chain with a PINN}

\subsection*{Case-study focus}
In this case study, you model a radioactive decay chain
\[
N_1 \rightarrow N_2 \rightarrow N_3
\]
with a direct supervised ANN and a physics-informed neural network (PINN). The system is governed by
\[
\frac{dN_1}{dt}=-\lambda_1N_1,\qquad
\frac{dN_2}{dt}=\lambda_1N_1-\lambda_2N_2,\qquad
\frac{dN_3}{dt}=\lambda_2N_2.
\]
The direct ANN learns from a small number of data points. The PINN additionally uses the differential-equation residuals as part of the loss. The central question is whether encoding physical structure in the learning objective improves the model's interpolation, extrapolation, and physical plausibility.

\subsection*{Learning goals}
After completing this activity, you should be able to:
\begin{itemize}
    \item Describe the decay-chain equations as a physics-based model of coupled time evolution.
    \item Explain how a direct ANN represents a mapping \(t\mapsto (N_1,N_2,N_3)\).
    \item Explain how a PINN adds physical constraints through differential-equation residuals.
    \item Compare data-driven fitting with physics-informed fitting under sparse-data conditions.
    \item Evaluate model outputs using both numerical error and physical criteria such as residual size, extrapolation behavior, non-negativity, and total-population consistency.
    \item Identify limitations of PINNs, including sensitivity to loss balancing, collocation-point choice, optimizer behavior, and the difference between satisfying a residual approximately and understanding the physics.
\end{itemize}

\subsection*{Before running the notebook}
Answer:
\begin{enumerate}
    \item What physical process is represented by each term in the three differential equations?
    \item What should happen to \(N_1(t)\), \(N_2(t)\), and \(N_3(t)\) over time?
    \item Is the total population \(N_1+N_2+N_3\) conserved in this closed decay chain? Show this from the equations.
    \item Why might a direct ANN trained on only a few points extrapolate poorly?
\end{enumerate}

\subsection*{Notebook workflow}

\paragraph{Step 1: Simulate reference data.}
Run the simulation cell. The notebook uses
\[
\lambda_1=0.5,\qquad \lambda_2=0.3,
\]
with initial values \(N_1(0)=0.8\), \(N_2(0)=0.1\), and \(N_3(0)=0.5\). The reference solution is generated by numerical Euler integration over the training interval.

Record:
\begin{itemize}
    \item training time interval,
    \item extrapolation interval,
    \item number of collocation points,
    \item initial total population \(N_1(0)+N_2(0)+N_3(0)\).
\end{itemize}

\paragraph{Step 2: Inspect the ANN architecture.}
The model maps one input, time \(t\), to three outputs, \((N_1,N_2,N_3)\). The hidden layers use \(\tanh\) activations.

Answer:
\begin{enumerate}
    \item Why is the output dimension three?
    \item What does it mean to treat the ANN as a parameterized function \(f_{\bm{\theta}}(t)\)?
    \item Does the architecture itself enforce the decay equations, non-negativity, or conservation of total population?
\end{enumerate}

\paragraph{Step 3: Train the direct ANN.}
Run the direct ANN section. This model is trained only on the selected collocation data values.

Answer:
\begin{enumerate}
    \item How many data points are used for training?
    \item Does a low training loss at the collocation points guarantee accurate behavior between the points?
    \item What kind of information is missing from the direct ANN loss?
\end{enumerate}

\paragraph{Step 4: Train the PINN.}
Run the PINN section. The PINN loss contains two parts:
\[
\mathcal{L}_{\mathrm{PINN}}
=
\mathcal{L}_{\mathrm{data}}
+
\mathcal{L}_{\mathrm{ODE}},
\]
where \(\mathcal{L}_{\mathrm{ODE}}\) penalizes violations of the decay-chain equations.

Answer:
\begin{enumerate}
    \item How are automatic derivatives used to compute \(dN_i/dt\)?
    \item What does each residual term \(r_1,r_2,r_3\) measure physically?
    \item Why can an ODE residual act as a learning bias?
    \item Does the PINN use physics as data, as architecture, or as a loss constraint?
\end{enumerate}

\paragraph{Step 5: Compare interpolation and extrapolation.}
Run the evaluation and plotting cells. The notebook compares the direct ANN and the PINN both within the training interval and in an extrapolation interval.

Complete:
\begin{center}
\begin{tabular}{l c c}
\hline
Diagnostic & Direct ANN & PINN \\
\hline
Training-region MSE & & \\
ODE residual loss & -- & \\
Extrapolation error & & \\
Qualitative extrapolation behavior & & \\
\hline
\end{tabular}
\end{center}

Answer:
\begin{enumerate}
    \item Which model follows the reference solution better in the extrapolation region?
    \item Does the PINN improvement come from more data, more physics, or both?
    \item Are there time regions where either model gives physically implausible predictions?
\end{enumerate}

\paragraph{Step 6: Add physical consistency checks.}
Use the plotted predictions to check:
\begin{itemize}
    \item whether \(N_1,N_2,N_3\) remain non-negative,
    \item whether \(N_1+N_2+N_3\) remains approximately constant,
    \item whether \(N_1\) decreases monotonically,
    \item whether \(N_3\) increases monotonically in the modeled interval.
\end{itemize}

These checks are not merely numerical. They test whether the learned function respects qualitative physics.

\subsection*{Synthesis task}
Write a 300--500 word comparison of the direct ANN and the PINN. Your explanation should include:
\begin{itemize}
    \item what information each model receives during training,
    \item how the loss functions differ,
    \item which model generalizes better and why,
    \item one physical diagnostic beyond MSE,
    \item one limitation of the PINN approach.
\end{itemize}

\subsection*{Optional extension}
Vary the number and location of collocation points. Investigate whether the PINN still outperforms the direct ANN when the collocation points are clustered in one part of the time interval. Explain the result in terms of data coverage and physics-based constraints.
